% from ega1-3-fr.tex
\oldpage[I]{115}
Comme $f_{(S')}=f\times_S 1_{S'}$, on voit que si pour tout préschéma $X$
l'\emph{identité} $1_X$ possède la propriété $\mathbf{P}$, (i) implique
(ii)~; comme d'autre part, $f\times_S g$ est le morphisme composé
\phantomsection\label{I.3.5.1.product-composition-fr}
\[
  X\times_S Y\xrightarrow{f\times 1_Y}X'\times_S Y
  \xrightarrow{1_{X'}\times g}X'\times_S Y'
\]
on voit que, si le \emph{composé} de deux morphismes possédant la propriété
$\mathbf{P}$, possède aussi cette propriété, alors (ii) implique (i).
\end{env}

Une première application de cette remarque est la

\begin{proposition}[3.5.2]
\label{I.3.5.2-fr}
\begin{enumerate}
  \item[(i)] Si $f:X\to X'$, $g:Y\to Y'$ sont des $S$-morphismes surjectifs,
  $f\times_S g$ est surjectif.
  \item[(ii)] Si $f:X\to Y$ est un $S$-morphisme surjectif, $f_{(S')}$ est
  surjectif pour toute extension $S'$ du préschéma de base.
\end{enumerate}
\end{proposition}

Le composé de deux surjections étant une surjection, il suffit de démontrer
(ii)~; mais cette proposition résulte aussitôt de
(\hyperref[I.3.4.8-fr]{3.4.8}) appliqué à $M=X$.

\begin{proposition}[3.5.3]
\label{I.3.5.3-fr}
Pour qu'un morphisme $f:X\to Y$ soit surjectif, il faut et il suffit que pour
tout corps $K$ et tout morphisme $\operatorname{Spec}(K)\to Y$, il existe une
extension $K'$ de $K$ et un morphisme $\operatorname{Spec}(K')\to X$ rendant
commutatif le diagramme
\phantomsection\label{I.3.5.3.surjectivity-diagram-fr}
\[
  \xymatrix{
    X\ar[d]_f & \operatorname{Spec}(K')\ar[l]\ar[d]\\
    Y & \operatorname{Spec}(K)\ar[l]
  }
\]
\end{proposition}

La condition est suffisante, car pour tout $y\in Y$, il suffit de l'appliquer
à un morphisme $\operatorname{Spec}(K)\to Y$ correspondant à un monomorphisme
$k(y)\to K$, $K$ étant une extension de $k(y)$
(\hyperref[I.2.4.6-fr]{2.4.6}). Inversement, supposons $f$ surjectif, et soit
$y\in Y$ l'image de l'unique point de $\operatorname{Spec}(K)$~; il existe
$x\in X$ tel que $f(x)=y$~; considérons le monomorphisme correspondant
$k(y)\to k(x)$ (\hyperref[I.2.2.1-fr]{2.2.1})~; il suffit alors de prendre
$K'$ extension de $k(y)$ telle qu'il existe des $k(y)$-monomorphismes de
$k(x)$ et de $K$ dans $K'$ (Bourbaki, \emph{Alg.}, chap. V, \S~4, prop. 2)~;
le morphisme $\operatorname{Spec}(K')\to X$ correspondant à $k(x)\to K'$
répond à la question.

Avec le langage introduit dans (\hyperref[I.3.4.5-fr]{3.4.5}), on peut dire
encore que \emph{tout point géométrique de $Y$ à valeurs dans $K$ provient
d'un point géométrique de $X$ à valeurs dans une extension de $K$}.

\begin{definition}[3.5.4]
\label{I.3.5.4-fr}
On dit qu'un morphisme de préschémas $f:X\to Y$ est \emph{universellement
injectif}, ou un \emph{morphisme radiciel}, si pour tout corps $K$,
l'application correspondante $X(K)\to Y(K)$ est injective.
\end{definition}

Il résulte aussitôt des définitions que tout \emph{monomorphisme de
préschémas} (T, 1.1) est radiciel.

\begin{env}[3.5.5]
\label{I.3.5.5-fr}
Pour qu'un morphisme $f:X\to Y$ soit radiciel, il suffit que la condition de
la déf. (\hyperref[I.3.5.4-fr]{3.5.4}) soit vérifiée pour tout corps
\emph{algébriquement clos}. En effet, si $K$ est un corps quelconque, $K'$ une
extension algébriquement close de $K$, le diagramme
\phantomsection\label{I.3.5.5.algebraic-closure-diagram-fr}
\[
  \xymatrix{
    X(K)\ar[r]^\alpha\ar[d]_\varphi & Y(K)\ar[d]^{\varphi'}\\
    X(K')\ar[r]_{\alpha'} & Y(K')
  }
\]

